\section*{Chapter \ref{ch:nw:long-haul-fibre} --- Long-Haul Fibre}

\begin{solution}{exo:nw:long-haul-fibre:1}
$0.16 \times 80 = \qty{12.8}{\decibel}$, a factor $10^{-1.28}$: about 5.2\% of the light is left.
\end{solution}

\begin{solution}{exo:nw:long-haul-fibre:2}
$(6\,550 - 21.7) \times 10^{-6} \times c_0 / 1.462 = \qty{1339}{\kilo\metre}$ of glass, a path factor of 1.185 over the \qty{1129.2}{\kilo\metre}
geodesic.
\end{solution}

\begin{solution}{exo:nw:long-haul-fibre:3}
With \omterm{def:nw:long-haul-fibre:dark}{dark fibre} the buyer chooses and controls the equipment (transponders, FEC, \omterm{def:nw:long-haul-fibre:disp}{dispersion compensation}, no switches), so it controls every
microsecond beyond the glass; with a \omterm{def:nw:long-haul-fibre:dark}{wavelength service} the seller's equipment and its choices set the latency, and may change when the seller
upgrades.
\end{solution}

\begin{solution}{exo:nw:long-haul-fibre:4}
Optical compensation adds a length of dispersion-compensating fibre to every span, and light takes time to cross it: in the model a fifth of the
path, more than a millisecond on the Chicago route. Electronic compensation undoes the dispersion by computation in the receiver, which costs a
processing delay at the end of the route, independent of its length.
\end{solution}

\begin{solution}{exo:nw:long-haul-fibre:5}
Across the ocean a cable can lie close to the geodesic, with detours only at the landings and the tails to the data centres; on land the route
must follow the rights of way it can obtain (roads, railways, utility corridors) around terrain and property. The Atlantic route is also five
times longer, so fixed detours at its ends weigh less in proportion.
\end{solution}

\begin{solution}{exo:nw:long-haul-fibre:6}
Whether both directions use the same fibres, path and equipment (otherwise the halves differ, and a round trip cannot tell them apart), whether it
is rack to rack or between the seller's own equipment, and whether it is a typical, a best or a guaranteed figure.
\end{solution}

\begin{solution}{exo:nw:long-haul-fibre:7}
The microwave route at 1.011 in air takes $1.011 \times 3\,767.8 = \qty{3809}{\micro\second}$ between the two buildings. Hollow-core glass along
the geodesic takes \qty{3778}{\micro\second}; with the model's \qty{21.7}{\micro\second} of equipment the hollow-core path factor must be at most
$(3\,809 - 21.7)/3\,778 = 1.0025$ (1.008 without equipment). A hollow-core route would have to be as straight as the microwave route and
carry almost no equipment delay.
\end{solution}

\begin{solution}{exo:nw:long-haul-fibre:8}
A path 3\% longer than the geodesic in glass takes $1.03 \times 1.462 = 1.51$ times the vacuum floor; a microwave route at a \omterm{def:nw:the-north-american-map:factor}{route factor} of 1.01
to 1.10 in air takes 1.01 to 1.10 times it. The zig-zag of the towers costs a few per cent; the glass costs 46\%. Straightness cannot make up
for the medium.
\end{solution}

\begin{solution}{pb:nw:long-haul-fibre:1}
\textbf{Part I.}
\begin{enumerate}
\item \qty{1129.2}{\kilo\metre}; \qty{3.767}{\milli\second} in vacuum, \qty{5.507}{\milli\second} in fibre.
\item \qty{13.3}{\milli\second} round trip against a round-trip fibre floor of \qty{11.014}{\milli\second}.
\item $6\,650 - 5\,507 = \qty{1143}{\micro\second}$.
\item Because the route is glass: even along the geodesic it cannot beat $n_g d/c_0$; the vacuum floor measures what the medium costs, not what
  the route wastes.
\end{enumerate}
\textbf{Part II.}
\begin{enumerate}[resume]
\item \qty{1359}{\kilo\metre} of glass, in 17 spans of at most \qty{80}{\kilo\metre}.
\item \qty{21.7}{\micro\second} (two terminals and 17 amplifiers): 1.9\% of the excess.
\item \qty{1121}{\micro\second} of extra path, \qty{230}{\kilo\metre} of glass.
\item The path shrinks to \qty{1323}{\kilo\metre}, still \qty{193}{\kilo\metre} (\qty{943}{\micro\second}) more than the geodesic: most of the
  excess remains path.
\end{enumerate}
\textbf{Part III.}
\begin{enumerate}[resume]
\item \qty{1.33}{\milli\second}: \qty{7.976}{\milli\second} in all.
\item \qty{4.569}{\milli\second}, \qty{2.08}{\milli\second} less.
\item \qty{5.528}{\milli\second}: straightening would save \qty{1.12}{\milli\second}, hollow core \qty{2.08}{\milli\second}.
\item The microwave route reaches Carteret from Aurora, \qty{52}{\kilo\metre} further west, in \qty{3.982}{\milli\second}: \qty{0.59}{\milli\second}
  faster than hollow core on the 2010 path, \qty{2.67}{\milli\second} faster than the 2010 route.
\end{enumerate}
\textbf{Part IV.}
\begin{enumerate}[resume]
\item \emph{Named result}: of the \qty{13.3}{\milli\second} round trip, \qty{11.01}{\milli\second} is glass along the geodesic,
  \qty{2.24}{\milli\second} extra path and \qty{0.04}{\milli\second} equipment (with the model's assumptions); \omterm{def:nw:long-haul-fibre:hcf}{hollow-core fibre} on the same path would
  save \qty{2.08}{\milli\second} one way, \qty{4.16}{\milli\second} round trip.
\item The split between path and equipment: every microsecond of equipment assumed is \qty{205}{\metre} less path.
\item Microwave networks along nearly the same geodesic, in a medium 46\% faster, came within 1 to 2\% of the air floor.
\item Bandwidth and reliability: it carries market data and orders in any weather and backs up the radio links.
\item It depends on the strategy's life and on control: for a firm that must control every piece of equipment for years, \omterm{def:nw:long-haul-fibre:dark}{dark fibre}; for a
  back-up or a data route, a leased wavelength.
\item A path factor within a fraction of a per cent of the geodesic and almost no equipment delay (exercise~\ref{exo:nw:long-haul-fibre:7}).
\item That on land the time is lost to the path: the Atlantic cable, with no rights of way to follow, has a factor of 1.09 against the Chicago
  route's 1.20.
\item Eleven go to light in glass along the geodesic, two to the detours of the path, and almost nothing to the equipment.
\end{enumerate}
\end{solution}

\begin{solution}{iq:nw:long-haul-fibre:1}
At $c_0/1.46$, about \qty{205000}{\kilo\metre\per\second}: \qty{4.88}{\micro\second} per kilometre, \qty{4.88}{\milli\second} per
\qty{1000}{\kilo\metre} one way, against \qty{3.34}{\milli\second} in vacuum or air.
\iqlookfor{The \omterm{def:nw:the-north-american-map:floor}{group index}; one-way against round trip; the per-kilometre rule of thumb.}
\end{solution}

\begin{solution}{iq:nw:long-haul-fibre:2}
The path's detours from the geodesic, optical \omterm{def:nw:long-haul-fibre:disp}{dispersion compensation} (extra glass), amplifiers (small), terminals: transponders and \omterm{def:nw:long-haul-fibre:fec}{forward
error correction}, and any switches or routers on the way; plus the tails from the route's ends to the racks.
\iqlookfor{Path length dominates; DCF against electronic compensation; FEC as a trade of distance against delay.}
\end{solution}

\begin{solution}{iq:nw:long-haul-fibre:3}
Microwave travels in air at almost $c_0$, 46\% faster than light in glass, along a nearly straight line of towers; fibre follows rights of way
and is slower per kilometre. Fibre remains for bandwidth, reliability in bad weather, and everything that is not the few latency-critical
messages.
\iqlookfor{Medium and path; capacity and weather; the two used together.}
\end{solution}

\begin{solution}{iq:nw:long-haul-fibre:4}
A long-term, usually prepaid, exclusive right to use specified fibres or capacity without owning the cable. A firm buys \omterm{def:nw:long-haul-fibre:dark}{dark fibre} when it needs
to control the equipment and hence the latency, expects to use the route for longer than the break-even term, and has the staff to run it;
otherwise it leases.
\iqlookfor{Control of latency; up-front against monthly cost; maintenance obligations; strategy lifetime.}
\end{solution}

\begin{solution}{iq:nw:long-haul-fibre:5}
Fibre that guides light in an air core, so light travels at nearly $c_0$; recent designs also lose less light than solid glass. For trading it
would cut about 30\% from fibre latency, bringing fibre routes close to the air floor with fibre's bandwidth and weather-independence, and
threatening microwave's advantage on routes that can be built straight.
\iqlookfor{Latency and loss both; the route must still be straight; deployment and cost as open questions.}
\end{solution}

\begin{solution}{iq:nw:long-haul-fibre:6}
The endpoints differ from the ones you computed (a nearer building); the figure is not one-way (a round trip halved over asymmetric paths, or
a wrong unit); the route is not all standard fibre (hollow core, or part radio); or the measurement is wrong (unsynchronised clocks). Ask
for the endpoints, the method and the medium, and measure.
\iqlookfor{Physics as a sanity check; endpoints and definitions; measurement.}
\end{solution}
